% discussion.tex -- EVM Executor Paper (BNR-2026-008)

\section{Discussion}

\subsection{Key Findings}

\textbf{Finding 1: Execution and proving must be architecturally separated.}
All five surveyed production systems maintain a clear boundary between the
execution engine (which processes transactions and generates traces) and
the proving engine (which consumes traces and generates ZK proofs). This
separation is not merely an implementation choice but an architectural
necessity: the EVM's 256-bit arithmetic, stack-based execution model, and
Keccak-centric hashing are fundamentally mismatched with the finite field
arithmetic and polynomial constraints of ZK proof systems. Attempting to
prove execution inline---without a trace intermediary---would require
encoding the full EVM interpreter in a constraint system, an approach no
production system has found practical.

\textbf{Finding 2: Import-as-module (Strategy~A) is viable for Basis Network.}
Optimism's op-geth demonstrates that a complete L2 execution layer can be
built with only 34~lines of EVM modification. The \texttt{vm.StateDB}
interface provides sufficient abstraction to integrate a custom Poseidon-based
sparse Merkle tree without modifying Geth's internal state management.
Scroll's experience with zktrie (and subsequent migration away from it in the
Euclid upgrade) validates that tight coupling between the state trie and the
EVM fork creates maintenance burden that outweighs the benefits.

\textbf{Finding 3: Geth's tracing hooks API enables efficient ZK trace capture.}
The \texttt{core/tracing.Hooks} interface, available since Geth~1.14+, provides
event-driven callbacks for all state-modifying operations. By capturing only
state changes (10--30\% overhead) rather than every opcode execution
(10--50$\times$ overhead with structLogger), the ZK tracer achieves throughput
compatible with enterprise requirements.

\textbf{Finding 4: KECCAK256 is the dominant ZK constraint cost.}
At ${\sim}$150{,}000 R1CS constraints per invocation (approximately 1{,}000$\times$
the cost of Poseidon), KECCAK256 represents the single largest proving cost
in any zkEVM. Two mitigation strategies are available: (a)~replacing Keccak
with Poseidon where possible (state trie, address derivation) and (b)~using
a preimage oracle with lookup tables for Keccak values that must be verified
(contract storage layout, ABI encoding).

\textbf{Finding 5: Scroll is moving away from custom zktrie.}
The 2025 Euclid upgrade migrates from the Poseidon-based zktrie back to
standard MPT combined with OpenVM (a general-purpose zkVM). This suggests
that the approach of modifying the state trie within the Geth fork---rather
than abstracting it behind an interface---creates unsustainable coupling.
Our Strategy~A approach (custom StateDB behind \texttt{vm.StateDB} interface)
avoids this pitfall.

\subsection{Invariants Discovered}

The following invariants are established by this investigation and must be
preserved by downstream pipeline stages:

\begin{enumerate}
  \item \textbf{INV-EVM-1 (Execution-Proving Separation)}: The EVM executor
    SHALL produce execution traces consumed asynchronously by the ZK prover.
    No inline proving within the execution loop.

  \item \textbf{INV-EVM-2 (Trace Completeness)}: Every state-modifying operation
    (SLOAD, SSTORE, balance change, nonce change, LOG, CREATE) SHALL appear in
    the execution trace. Omitting any operation invalidates the ZK witness.

  \item \textbf{INV-EVM-3 (Cancun Compatibility)}: The execution engine SHALL
    support all Cancun EVM opcodes. No Pectra-era opcodes (PragueTime must
    remain unset) due to the Avalanche Subnet-EVM constraint.

  \item \textbf{INV-EVM-4 (KECCAK256 Mitigation)}: KECCAK256 operations SHALL
    be handled via preimage oracle with lookup tables in the ZK circuit,
    not proved directly in R1CS or PLONKish constraints.

  \item \textbf{INV-EVM-5 (Module Import Strategy)}: The execution engine
    SHALL import go-ethereum as a Go module dependency (Strategy~A) unless
    the \texttt{vm.StateDB} interface proves insufficient for Poseidon SMT
    integration, in which case a minimal fork (Strategy~B) is acceptable.

  \item \textbf{INV-EVM-6 (Selective Tracing)}: The ZK tracer SHALL capture
    state-modifying operations only (not every opcode) in production mode.
    Full opcode capture SHALL be available as a debug mode only.
\end{enumerate}

\subsection{Precompile Considerations}

EVM precompiled contracts (addresses 0x01--0x0a) require special handling
in the ZK circuit:

\begin{table}[htbp]
\caption{Precompile ZK Constraint Estimates}
\label{tab:precompiles}
\centering
\begin{tabular}{llrl}
\toprule
\textbf{Address} & \textbf{Name} & \textbf{Constraints} & \textbf{Notes} \\
\midrule
0x01 & ecRecover & ${\sim}$30{,}000 & secp256k1 recovery \\
0x02 & SHA256 & ${\sim}$30{,}000 & Boolean circuit \\
0x03 & RIPEMD160 & ${\sim}$30{,}000 & Boolean circuit \\
0x04 & identity & Trivial & Data copy \\
0x05 & modexp & Variable & Dynamic cost \\
0x06 & ecAdd (BN254) & ${\sim}$1{,}000 & Native to ZK field \\
0x07 & ecMul (BN254) & ${\sim}$5{,}000 & Native to ZK field \\
0x08 & ecPairing & ${\sim}$100{,}000 & Per pair element \\
0x09 & blake2f & ${\sim}$10{,}000 & Compression function \\
0x0a & KZG eval & ${\sim}$50{,}000 & Cancun (EIP-4844) \\
\bottomrule
\end{tabular}
\end{table}

BN254 operations (ecAdd, ecMul) are natively efficient because the ZK proof
system itself operates over the BN254 curve. KECCAK256-derived precompiles
(ecRecover via message hash, SHA256, RIPEMD160) inherit the high constraint
cost of Boolean circuit emulation.

\subsection{Limitations}

\textbf{Projected benchmarks.} The throughput figures in this paper are
projections derived from published data, not direct measurements. While
cross-validated against multiple independent sources, actual performance
on Basis Network hardware may differ. The Go benchmark code is provided
for execution once the Go runtime is available.

\textbf{Trace format stability.} The ZK trace format defined in this paper
is preliminary. The downstream witness generation module (RU-L3) will
specify the exact trace schema consumed by the Rust prover, which may
require modifications to the tracer output structure.

\textbf{StateDB interface sufficiency.} The recommendation to use Strategy~A
(import as module) rests on the assumption that \texttt{vm.StateDB} can
accommodate a Poseidon SMT. Scroll's zktrie experience validates this
assumption, but edge cases (state root computation during EVM execution,
snapshot/revert semantics) may require interface extensions. If extensions
prove necessary, a minimal fork targeting only \texttt{core/state} becomes
the fallback.

\textbf{Single-threaded execution.} All projections assume single-threaded
EVM execution. Parallel execution (as demonstrated by Gravity Reth at
41{,}000~tx/s for ERC20 transfers) could improve throughput substantially
but introduces complexity in trace ordering and state consistency.
